\chapter*{Literature and Context}
\addcontentsline{toc}{chapter}{Literature and Context}
\markboth{LITERATURE AND CONTEXT}{LITERATURE AND CONTEXT}

This bibliography is deliberately annotated. The building blocks
of Horimetrik -- mixed-radix systems, falling factorials,
inversion sequences, the calculus of finite differences -- are
established mathematics, and this book claims nothing else. Each
entry therefore states two things: which building block it
contributes, and what is \emph{not} to be found there -- that is,
the point at which Horimetrik goes beyond its source. Whoever
wants to check the placement given in
Chapter~\ref{ch:einordnung} will find the evidence here.

% thebibliography would otherwise set its own \chapter*
% ("Bibliography") on a new page; we already have a chapter
% heading above.
\begingroup
\makeatletter
\renewcommand{\chapter}[2]{}%
\renewcommand{\@mkboth}[2]{}%
\makeatother
\begin{thebibliography}{10}

\bibitem{cantor} G.~Cantor: \emph{Über die einfachen Zahlensysteme},
Zeitschrift für Mathematik und Physik 14 (1869), 121--128.
--- The factorial base for fractions; this is exactly the
fraction-preserving limit space of the third staircase
(Chapter~\ref{ch:erweiterungen}), but not the shape-preserving
world, nor the interplay between them.

\bibitem{knuth} D.~E. Knuth: \emph{The Art of Computer Programming},
vol.~2, 3rd ed., Addison--Wesley 1997, Section 4.1.
--- Mixed-radix systems and the factorial number system
(factoradic); there the bases grow towards the significant side,
and leading zeros are value-neutral.

\bibitem{lehmer} D.~H. Lehmer: \emph{Teaching combinatorial tricks
to a computer}, Proc. Sympos. Appl. Math. 10 (1960), 179--193.
--- Lehmer codes and inversion tables; they explain the factorial
size of the horizons and the digit bound $0\le a_i\le i$, but not
the horizon-dependent evaluation.

\bibitem{gkp} R.~L. Graham, D.~E. Knuth, O.~Patashnik:
\emph{Concrete Mathematics}, 2nd ed., Addison--Wesley 1994, Ch.~2.
--- Falling factorials and the calculus of finite differences; the
toolkit of the horizon derivative.

\bibitem{cahen} P.-J. Cahen, J.-L. Chabert: \emph{Integer-Valued
Polynomials}, AMS Surveys and Monographs 48, 1997.
--- Polynomials with integer values in the binomial basis;
Horimetrik uses only the positive cone and evaluates at the
running horizon.

\bibitem{rigo} M.~Rigo, M.~Stipulanti: \emph{Revisiting regular
sequences in light of rational base numeration systems},
arXiv:2103.16966.
--- Abstract numeration systems in which leading zeros are
value-sensitive via the radix order: the closest known neighbour
of the zero semantics, treated there in formal-language terms
rather than algebraically.

\bibitem{kudr} G.~Kudryavtseva, V.~Mazorchuk: \emph{On Kiselman's
semigroup}, arXiv:math/0511374; O.~Ganyushkin, V.~Mazorchuk:
\emph{On Kiselman quotients of 0-Hecke monoids}, arXiv:1006.0316.
--- Monoids of idempotent generators with absorption relations;
the neighbourhood of the six-element compactification monoid.

\bibitem{simon} D.~Simon: \emph{Mixed radix numeration bases:
Horner's rule, Yang-Baxter equation and Furstenberg's conjecture},
arXiv:2405.19798.
--- Local base changes in mixed-radix systems; a candidate tool
for the interchange defect.

\bibitem{invseq} S.~Corteel, M.~A.~Martinez, C.~D.~Savage,
M.~Weselcouch: \emph{Patterns in Inversion Sequences I},
arXiv:1510.05434.
--- Inversion sequences ($0\le e_j<j$) and their bijection with
permutations; our $H_n$ is, as a set, exactly this space
(Chapter~\ref{ch:einordnung}). What is not there: the reversed
valuation, the horizon transitions, and the structure horizon as a
statistic.

\bibitem{lecturehall} C.~D. Savage, M.~J. Schuster: \emph{Ehrhart
series of lecture hall polytopes and Eulerian polynomials for
inversion sequences}, J. Combin. Theory Ser.~A 119 (2012),
850--870.
--- $s$-inversion sequences with arbitrary capacity profiles and
their geometry; the framework for the $s$-Horimetrik programme of
Chapter~\ref{ch:einordnung}.

\bibitem{tightentries} S.~Chern, S.~Fu, Z.~Lin: \emph{Burstein's
permutation conjecture, Hong and Li's inversion sequence
conjecture, and restricted Eulerian distributions},
arXiv:2209.12137.
--- Statistics on inversion sequences, among them the \emph{tight
entries} ($e_i=s_i-1$, statistic $\mathrm{tig}$); the structure
horizon refines this extreme-case count into a distance measure
(Chapter~\ref{ch:einordnung}).

\bibitem{middleorder} M.~Bouvel, L.~Ferrari, B.~E.~Tenner:
\emph{Between weak and Bruhat: the middle order on permutations},
arXiv:2405.08943.
--- The coordinatewise order on inversion codes as a lattice
between the weak and the Bruhat order; the framework in which
$\mathcal F_r$ becomes a principal ideal and $\sigma$ a level of
ideals.

\bibitem{umbral} K.~Beauduin: \emph{Operational Umbral Calculus},
arXiv:2407.16348.
--- Modern finite operator calculus: delta operators and their
basic sequences. The framework for $\Delta$ and the structure
polynomials; what is not there is the capacity bookkeeping
$\Gamma_T$.

\bibitem{ikenmeyerpak} C.~Ikenmeyer, I.~Pak: \emph{What is in \#P
and what is not?}, FOCS 2022, arXiv:2204.13149.
--- Characterizes the \emph{binomial-good} polynomials
(nonnegative integer coefficients in the binomial basis) as the
\emph{relativizing} polynomial closure operations of \#P, and
states the binomial basis conjecture (without relativization the
same characterization is only conjectural), whose univariate
version would already imply P$\,\ne\,$NP. Because of
$\ff Xd=d!\binom Xd$ the shapes of Horimetrik are a proper
subclass of the binomial-good polynomials
(Chapter~\ref{ch:strukturpolynome}); what is \emph{not} there is
the quantitative layer: the height $\max(d+c_d)$, the $(r{+}1)!$
hierarchies, and their growth laws.

\bibitem{oeis} OEIS Foundation Inc.: \emph{The On-Line Encyclopedia
of Integer Sequences}, \texttt{https://oeis.org}.
--- Reference database for the nine counting sequences of this
book; A-numbers will be added once assigned.

\end{thebibliography}
\endgroup
